\RequirePackage[l2tabu, orthodox]{nag}
\pdfinfoomitdate=1
\pdfsuppressptexinfo=15
\pdftrailerid{}
\providecommand{\BotgeDarkMode}{1}
\documentclass[11pt,a4paper]{article}

\usepackage[T1]{fontenc}
\usepackage{lmodern}
\usepackage{microtype}
\usepackage[margin=1in]{geometry}
\usepackage{amsmath,amssymb,amsthm,mathtools}
\usepackage{booktabs}
\usepackage{xcolor}

\definecolor{BotgeBg}{HTML}{0e1117}
\definecolor{BotgeFg}{HTML}{e9edf4}
\definecolor{BotgeAccent}{HTML}{58a6ff}
\definecolor{BotgeLightAccent}{HTML}{0969da}

\usepackage{darkmode}
\ifnum\BotgeDarkMode=1\relax
  \enabledarkmode
  \pagecolor{BotgeBg}
  \color{BotgeFg}
  \colorlet{fgcolor}{BotgeFg}
\else
  \disabledarkmode
\fi

\usepackage{hyperref}
\hypersetup{%
  colorlinks=true,%
  bookmarksnumbered=true,%
  bookmarksopen=true,%
  pdfdisplaydoctitle=true,%
  pdfstartview=FitH,%
  pdftitle={A Deterministic Bimolecular Reaction Atlas and Bounded Catalog Projection},%
  pdfauthor={The Botge Project},%
  pdfsubject={Stoichiometric Conservation, Order Invariance, and Reactant-Avoiding Projection},%
  pdfkeywords={bimolecular organic reactions, SMARTS, InChIKey, stoichiometry, FNV-1a projection, formal verification}%
}
\IfDarkModeTF{%
  \hypersetup{linkcolor=BotgeAccent, citecolor=BotgeAccent, urlcolor=BotgeAccent}%
}{%
  \hypersetup{linkcolor=BotgeLightAccent, citecolor=BotgeLightAccent, urlcolor=BotgeLightAccent}%
}
\usepackage{aliascnt}
\usepackage[capitalise,noabbrev,nameinlink]{cleveref}

\newcommand{\SharedTheorem}[3]{%
  \newaliascnt{#1}{theorem}%
  \newtheorem{#1}[#1]{#2}%
  \aliascntresetthe{#1}%
  \crefname{#1}{#2}{#3}%
}
\theoremstyle{plain}
\newtheorem{theorem}{Theorem}[section]
\crefname{theorem}{Theorem}{Theorems}
\SharedTheorem{proposition}{Proposition}{Propositions}
\theoremstyle{definition}
\SharedTheorem{definition}{Definition}{Definitions}
\theoremstyle{remark}
\newtheorem*{remark*}{Remark}

\newcommand{\N}{\mathbb{N}}
\newcommand{\Z}{\mathbb{Z}}
\newcommand{\Atlas}{\mathcal{M}}
\newcommand{\Reacts}{\mathcal{R}}
\newcommand{\Conds}{\mathcal{C}}
\newcommand{\SNtwo}{S_{\mathrm{N}}2}
\newcommand{\FNV}{\operatorname{FNV1a}_{64}}
\newcommand{\mset}[1]{\lbrace#1\rbrace}
\newcommand{\chem}[1]{\ensuremath{\mathrm{#1}}}
\newcommand{\smiles}[1]{\texttt{#1}}

\title{%
  A Deterministic Bimolecular Reaction Atlas\\
  and Bounded Catalog Projection%
}
\author{The Botge Project}
\date{}

\begin{document}
\maketitle

\begin{abstract}
  We specify a deterministic, rule-based bimolecular organic reaction system defined over a
  curated atlas \( \Atlas \) of \( K = 32 \) neutral substrates, eight priority-ordered SMARTS
  reaction transforms \( \Reacts \), and four explicit reagent regimes \( \Conds \) paired with an
  automatic selector \( \Conds_* \). Every reactive transform conserves all non-hydrogen heavy
  atoms modulo its declared stoichiometric byproduct, passes valence sanitization to a unique
  canonical isomeric SMILES product invariant under reactant permutation, and maps the resulting
  standard \(\text{InChIKey}\) into bounded base and overlay index sets \( \Z_{M_b} \) and
  \( \Z_{M_o} \) (\( 2 \le M_b \le 1024 \), \( 2 \le M_o \le 128 \)) while strictly avoiding
  both reactant base indices whenever \( M_b \ge 3 \). Structural invariants are mechanized in
  Lean~4 with Mathlib, cross-checked symbolically in SymPy and Z3, and certified exhaustively
  across all \( 32 \times 32 \times 5 = 5{,}120 \) ordered reactant-regime triples.
\end{abstract}

\noindent\textbf{Keywords:} bimolecular organic transforms, SMARTS atom mapping, stoichiometry,
InChIKey, FNV-1a projection, order invariance, formal verification.

\section{Introduction}\label{sec:introduction}

Given two discrete catalog entries mapped onto a small molecular library and an operator-selected
reaction regime, we seek a deterministic synthesis rule whose outcome is chemically conservative,
independent of input ordering, and mapped back onto disjoint catalog indices without inventing
spurious products when two functional groups are mutually incompatible. We size the index pools and
substrate hash in~\cref{sec:pools}, define the eight bimolecular reaction templates and prove
heavy-atom stoichiometric conservation in~\cref{sec:reactions}, establish order invariance and
reactant-avoiding bounded projection in~\cref{sec:projection}, report exhaustive finite-atlas
certification in~\cref{sec:certificate}, and state the chemical scope and limitations of the model
in~\cref{sec:limitations}.

\paragraph{Design choices.} Rather than sampling arbitrary molecular pairs or allowing unconstrained
multi-component cascades, the system restricts every evaluation to two substrates drawn from a
fixed \( 32 \)-member organic atlas and eight classical \chem{C\text{-}C}, \chem{C\text{-}N}, and
\chem{C\text{-}O} bond-forming transforms executed under thermal, Br\o{}nsted acid, Br\o{}nsted base,
or palladium-catalyzed conditions~\cite{Baumann1886,DielsAlder1928,FischerSpeier1895,Schiff1864,
Schotten1884,Williamson1850,Michael1887,MiyauraSuzuki1995,Wurtz1872}. Leaving groups are
explicitly accounted for at each
template boundary, canonical product keys rely on standard InChIKey~\cite{Heller2015} and
canonical isomeric SMILES~\cite{Weininger1988}, and incompatible pairs terminate cleanly with
functional-group diagnostics rather than heuristic fallback substitutions.

\paragraph{Verification.} Pool bounds, FNV-1a index totality, reactant-order invariance,
bounded reactant-avoiding projection, and heavy-atom multiset conservation are formalized in
Lean~4~\cite{deMouraUllrich2021} using Mathlib~\cite{Mathlib2020} with only the three standard
Lean and Mathlib axioms (\texttt{propext}, \texttt{Classical.choice}, \texttt{Quot.sound}).
Symbolic linear integer conservation and threshold minimality are cross-checked with SymPy
and Z3, and exhaustive execution over all \( 5{,}120 \) ordered triples confirms zero valence
sanitization failures, zero stoichiometric mismatches, and exact commutativity.

\subsection{Notation}\label{sec:notation}

\begin{itemize}
  \item \textbf{Modular index sets.} For integers \( M \ge 2 \), write
    \( \Z_M \coloneqq \{0, \dots, M - 1\} \). For \( a, b \in \Z_K \), let
    \( \operatorname{canon}(a, b) \coloneqq (\min(a, b), \max(a, b)) \).
  \item \textbf{Molecular graphs and heavy-atom multisets.} A neutral organic molecule is modeled
    as a finite connected multigraph \( G = (V, E, \alpha, b, q, h) \), where each vertex
    \( v \in V \) carries a non-hydrogen atomic symbol \( \alpha(v) \in \Sigma \) over the fixed
    alphabet
    \[
      \Sigma \coloneqq \{\mathrm{B}, \mathrm{C}, \mathrm{N}, \mathrm{O}, \mathrm{Cl}, \mathrm{Br}\},
    \]
    rational bond orders \( b : E \to \{1, 1.5, 2, 3\} \), zero formal charge \( q(v) = 0 \), and
    implicit hydrogen count \( h(v) \in \N \) satisfying standard covalent valence and
    Kekul\'e-representation rules after
    sanitization. Define the heavy-atom element multiset \( \nu(G) : \Sigma \to \N \) by
    \[
      \nu(G)(Z) \coloneqq \lvert \{ v \in V \mid \alpha(v) = Z \} \rvert, \qquad Z \in \Sigma.
    \]
    Multiset addition \( + \) and subtraction \( - \) are pointwise on \( \N^{\Sigma} \).
  \item \textbf{Reaction nomenclature.} Eponymous binary names take en-dashes
    (\emph{Diels--Alder}, \emph{Suzuki--Miyaura}, \emph{Schotten--Baumann}), brackets
    \([4 + 2]\) denote ring-atom counts assembled into the cyclohexene product, and \( \SNtwo \)
    denotes concerted bimolecular nucleophilic substitution with italic \(S\) and upright subscript
    \(\mathrm{N}\).
\end{itemize}

\section{Bounded Pools and Curated Substrate Atlas}\label{sec:pools}

\begin{proposition}[Pool Budget]\label{prop:chem-budget}
  Let \( C_b, C_o \ge 2 \) be the available base and overlay catalog sizes, and define
  the capped pool sizes
  \begin{equation}\label{eq:chem-pool-cap}
    M_b = \min(C_b, 1024), \qquad M_o = \min(C_o, 128).
  \end{equation}
  Then \( 2 \le M_b \le 1024 \) and \( 2 \le M_o \le 128 \), so both finite index sets
  \( \Z_{M_b} \) and \( \Z_{M_o} \) are nonempty and have cardinality at least two.
  Whenever \( M_b \ge 3 \) and \( E \subseteq \Z_{M_b} \) satisfies \( \lvert E \rvert \le 2 \),
  the complement \( \Z_{M_b} \setminus E \) is nonempty.
\end{proposition}
\begin{proof}
  Since \( C_b \ge 2 \) and \( 1024 \ge 2 \), \( M_b = \min(C_b, 1024) \) satisfies
  \( 2 \le M_b \le 1024 \), and analogously \( 2 \le M_o \le 128 \). Hence
  \( \lvert \Z_{M_b} \rvert = M_b \ge 2 \) and \( \lvert \Z_{M_o} \rvert = M_o \ge 2 \). When
  \( M_b \ge 3 \) and \( \lvert E \rvert \le 2 \), cardinality monotonicity yields
  \( \lvert \Z_{M_b} \setminus E \rvert \ge M_b - \lvert E \rvert \ge 3 - 2 = 1 \), so
  \( \Z_{M_b} \setminus E \neq \emptyset \).
\end{proof}

\begin{definition}[Stable Substrate Map]\label{def:substrate-map}
  Let \( \FNV : {(\{0,\dots,255\})}^* \to \{0, \dots, 2^{64} - 1\} \) denote the 64-bit FNV-1a
  hash with offset basis \( \texttt{cbf29ce484222325}_{16} \) and prime
  \( \texttt{100000001b3}_{16} \).
  Each immutable identifier byte string \( s \) maps to a substrate index
  \begin{equation}\label{eq:substrate-index}
    \iota(s) = \FNV(s) \bmod K, \qquad K = 32,
  \end{equation}
  referencing the curated \( 32 \)-substrate organic atlas \( \Atlas \) in~\cref{tab:substrates}
  independently of all other entries in the catalog.
\end{definition}

\begin{table}[t]
  \centering
  \caption{Curated \( 32 \)-substrate organic atlas \( \Atlas \) with Hill molecular formulas and canonical SMILES.\label{tab:substrates}}
  \scriptsize
  \setlength{\tabcolsep}{4pt}
  \begin{tabular}{@{}rllll@{}}
    \toprule
    \(i\) & Substrate & Class & Formula & Canonical SMILES \\
    \midrule
    0 & 1,3-butadiene & conjugated diene & \chem{C_4H_6} & \smiles{C=CC=C} \\
    1 & isoprene & conjugated diene & \chem{C_5H_8} & \smiles{C=CC(=C)C} \\
    2 & 1,3-cyclohexadiene & cyclic diene & \chem{C_6H_8} & \smiles{C1=CCCC=C1} \\
    3 & cyclopentadiene & cyclic diene & \chem{C_5H_6} & \smiles{C1=CCC=C1} \\
    4 & maleic anhydride & cyclic dienophile & \chem{C_4H_2O_3} & \smiles{O=C1C=CC(=O)O1} \\
    5 & methyl acrylate & enoate / Michael acceptor & \chem{C_4H_6O_2} & \smiles{C=CC(=O)OC} \\
    6 & acrylonitrile & conjugated nitrile & \chem{C_3H_3N} & \smiles{C=CC\#N} \\
    7 & methyl vinyl ketone & enone / Michael acceptor & \chem{C_4H_6O} & \smiles{C=CC(C)=O} \\
    8 & acetic acid & carboxylic acid & \chem{C_2H_4O_2} & \smiles{CC(=O)O} \\
    9 & propanoic acid & carboxylic acid & \chem{C_3H_6O_2} & \smiles{CCC(=O)O} \\
    10 & benzoic acid & aryl carboxylic acid & \chem{C_7H_6O_2} & \smiles{O=C(O)c1ccccc1} \\
    11 & methanol & primary alcohol & \chem{CH_4O} & \smiles{CO} \\
    12 & ethanol & primary alcohol & \chem{C_2H_6O} & \smiles{CCO} \\
    13 & benzyl alcohol & benzylic alcohol & \chem{C_7H_8O} & \smiles{OCc1ccccc1} \\
    14 & methylamine & primary amine & \chem{CH_5N} & \smiles{CN} \\
    15 & ethylamine & primary amine & \chem{C_2H_7N} & \smiles{CCN} \\
    16 & aniline & arylamine & \chem{C_6H_7N} & \smiles{Nc1ccccc1} \\
    17 & pyrrolidine & secondary amine & \chem{C_4H_9N} & \smiles{C1CCNC1} \\
    18 & acetaldehyde & aldehyde & \chem{C_2H_4O} & \smiles{CC=O} \\
    19 & propanal & aldehyde & \chem{C_3H_6O} & \smiles{CCC=O} \\
    20 & benzaldehyde & aryl aldehyde & \chem{C_7H_6O} & \smiles{O=Cc1ccccc1} \\
    21 & acetone & ketone & \chem{C_3H_6O} & \smiles{CC(C)=O} \\
    22 & acetyl chloride & acyl chloride & \chem{C_2H_3ClO} & \smiles{CC(=O)Cl} \\
    23 & benzoyl chloride & aroyl chloride & \chem{C_7H_5ClO} & \smiles{O=C(Cl)c1ccccc1} \\
    24 & 1-bromobutane & primary alkyl bromide & \chem{C_4H_9Br} & \smiles{CCCCBr} \\
    25 & benzyl bromide & benzylic bromide & \chem{C_7H_7Br} & \smiles{BrCc1ccccc1} \\
    26 & bromobenzene & aryl bromide & \chem{C_6H_5Br} & \smiles{Brc1ccccc1} \\
    27 & 4-bromoanisole & aryl bromide & \chem{C_7H_7BrO} & \smiles{COc1ccc(Br)cc1} \\
    28 & phenylboronic acid & arylboronic acid & \chem{C_6H_7BO_2} & \smiles{OB(O)c1ccccc1} \\
    29 & 4-tolylboronic acid & arylboronic acid & \chem{C_7H_9BO_2} & \smiles{Cc1ccc(B(O)O)cc1} \\
    30 & acetylacetone & 1,3-diketone & \chem{C_5H_8O_2} & \smiles{CC(=O)CC(C)=O} \\
    31 & dimethyl malonate & 1,3-diester & \chem{C_5H_8O_4} & \smiles{COC(=O)CC(=O)OC} \\
    \bottomrule
  \end{tabular}
\end{table}

\section{Bimolecular Reactions and Stoichiometric Conservation}\label{sec:reactions}

Let \( \Conds = \{\text{heat}, \text{acid}, \text{base}, \text{palladium}\} \) denote the four
mutually exclusive explicit reagent regimes, corresponding chemically to thermal activation
(\( \Delta \)), Br\o{}nsted acid catalysis (\( \mathrm{H^+} \)), Br\o{}nsted base promotion
(\( \mathrm{B^-} \)), and palladium-mediated basic cross-coupling (\( \mathrm{Pd^0}/\text{base} \)).
Let \( \Conds_* = \Conds \cup \{\text{auto}\} \), where \(\text{auto}\) admits all transforms
across \( \Conds \) in strict priority order.

\begin{definition}[Curated Bimolecular Reaction Set]\label{def:reaction-set}
  The reaction family \( \Reacts = \{r_1, \dots, r_8\} \) consists of eight priority-ordered,
  bimolecular constitutional transforms summarized in~\cref{tab:templates}:
  \( r_1 \) Diels--Alder \([4 + 2]\) cycloaddition (\(\text{heat}\), no leaving group,
  \( \beta_1 = 0 \));
  \( r_2 \) Fischer esterification (\(\text{acid}\), leaving species \chem{H_2O},
  \( \beta_2 = \mset{\mathrm{O}} \));
  \( r_3 \) Imine condensation (\(\text{acid}\), leaving species \chem{H_2O},
  \( \beta_3 = \mset{\mathrm{O}} \));
  \( r_4 \) Acylation / amide coupling (\(\text{base}\), leaving species \chem{HCl},
  \( \beta_4 = \mset{\mathrm{Cl}} \));
  \( r_5 \) Williamson \( \SNtwo \) substitution (\(\text{base}\), leaving species \chem{HBr},
  \( \beta_5 = \mset{\mathrm{Br}} \));
  \( r_6 \) Aldol condensation (\(\text{base}\), leaving species \chem{H_2O},
  \( \beta_6 = \mset{\mathrm{O}} \));
  \( r_7 \) Michael conjugate addition (\(\text{base}\), 100\% atom economy, \( \beta_7 = 0 \));
  \( r_8 \) Suzuki--Miyaura cross-coupling (\(\text{palladium}\), leaving species
  \chem{B(OH)_2Br}, \( \beta_8 = \mset{\mathrm{B}, 2\mathrm{O}, \mathrm{Br}} \)).
\end{definition}

\begin{table}[t]
  \centering
  \caption{Priority-ordered bimolecular reaction templates \( r_1 \dots r_8 \), balanced schemes, and heavy-atom leaving multisets \( \beta_k \).\label{tab:templates}}
  \scriptsize
  \setlength{\tabcolsep}{4pt}
  \begin{tabular}{@{}rllll@{}}
    \toprule
    \(k\) & Reaction & Balanced scheme & Byproduct & \( \beta_k \) \\
    \midrule
    1 & Diels--Alder \([4 + 2]\) & \(\text{diene} + \text{dienophile} \xrightarrow{\Delta} \text{cyclohexene adduct}\) & none & \(0\) \\
    2 & Fischer esterification & \(\mathrm{RCO_2H + R'OH} \xrightarrow{\mathrm{H^+}} \mathrm{RCO_2R' + H_2O}\) & \chem{H_2O} & \( \mset{\mathrm{O}} \) \\
    3 & Imine condensation & \(\mathrm{RR'C{=}O + R''NH_2} \xrightarrow{\mathrm{H^+}} \mathrm{RR'C{=}NR'' + H_2O}\) & \chem{H_2O} & \( \mset{\mathrm{O}} \) \\
    4 & Acylation / amide & \(\mathrm{RCOCl + R'R''NH} \xrightarrow{\mathrm{B^-}} \mathrm{RCONR'R'' + HCl}\) & \chem{HCl} & \( \mset{\mathrm{Cl}} \) \\
    5 & Williamson \( \SNtwo \) & \(\mathrm{ROH + R'CH_2Br} \xrightarrow{\mathrm{B^-}} \mathrm{ROCH_2R' + HBr}\) & \chem{HBr} & \( \mset{\mathrm{Br}} \) \\
    6 & Aldol condensation & \(\mathrm{RCOCH_2R' + R''CHO} \xrightarrow{\mathrm{B^-}} \mathrm{RCOC(R'){=}CHR'' + H_2O}\) & \chem{H_2O} & \( \mset{\mathrm{O}} \) \\
    7 & Michael addition & \(\beta\text{-dicarbonyl} + \text{acceptor} \xrightarrow{\mathrm{B^-}} \text{1,5-adduct}\) & none & \(0\) \\
    8 & Suzuki--Miyaura & \(\mathrm{Ar\text{-}Br + Ar'\text{-}B{(OH)}_2} \xrightarrow{\mathrm{Pd^0}/\text{base}} \mathrm{Ar\text{-}Ar' + B{(OH)}_2Br}\) & \chem{B(OH)_2Br} & \( \mset{\mathrm{B}, 2\mathrm{O}, \mathrm{Br}} \) \\
    \bottomrule
  \end{tabular}
\end{table}

\begin{theorem}[Stoichiometric Heavy-Atom Conservation]\label{thm:stoichiometry}
  For every template \( r_k \in \Reacts \) with heavy-atom leaving multiset \( \beta_k \)
  and every matched substrate pair \( (A, B) \in \Atlas^2 \) yielding sanitized product \( P \),
  \begin{equation}\label{eq:stoichiometry}
    \nu(A) + \nu(B) = \nu(P) + \beta_k.
  \end{equation}
  In particular, \( \nu(P) \) is determined elementwise by subtracting \( \beta_k \)
  from \( \nu(A) + \nu(B) \), with zero net heavy-atom loss or creation.
\end{theorem}
\begin{proof}
  Each template \( r_k \in \Reacts \) defines a deterministic SMARTS reaction transform on an
  ordered pair of neutral connected molecular graphs \( (A, B) \) in which: (1)~every heavy
  atom appearing in the reactant pattern outside the leaving group carries a bijective atom map
  label preserved into the product pattern; (2)~every unmapped heavy atom in the reactant pattern
  belongs to an explicitly excised leaving group whose element multiset equals \( \beta_k \)
  (\( \mset{\mathrm{O}} \) for hydroxyl in \( r_2 \) and carbonyl oxygen in \( r_3, r_6 \);
  \( \mset{\mathrm{Cl}} \) for acyl chloride in \( r_4 \); \( \mset{\mathrm{Br}} \) for alkyl or
  aryl bromide in \( r_5, r_8 \); and \( \mset{\mathrm{B}, 2\mathrm{O}} \) for the boronic acid
  group in \( r_8 \)); and (3)~unmapped spectator atoms outside the reaction center retain their
  element identities identically from \( A \) and \( B \) into the single connected product graph
  \( P \). Because graph valence sanitization only recomputes aromaticity, hybridization, and
  implicit hydrogen counts without adding or deleting heavy vertices, summing vertex counts by
  element over \( \Sigma \) yields \( \nu(P) = \nu(A) + \nu(B) - \beta_k \),
  establishing~\eqref{eq:stoichiometry} whenever a sanitized unimolecular product \( P \) matches a
  single transform application; an enforced multiset guard verifies~\eqref{eq:stoichiometry}
  dynamically before admitting any candidate.
\end{proof}

\begin{remark*}
  Hydrogen atoms are implicit in the graph vertex multiset \( \nu(G) \), and the \textsf{Byproduct}
  column in~\cref{tab:templates} lists formal neutral stoichiometric equivalents
  (\chem{H_2O}, \chem{HCl}, \chem{HBr}, \chem{B(OH)_2Br}) for heavy-atom mass balance; under
  basic conditions, halide acids and bromoboronic acid neutralize to halide and borate salts. Two
  additions (\( r_1, r_7 \)) have \( \beta_k = 0 \) and exhibit \( 100\% \) atom economy, whereas
  all six condensations, substitutions, and couplings satisfy exact mass balance upon inclusion of
  their formal byproducts. Template \( r_6 \) covers aldehyde, ketone, enone, and Knoevenagel
  \( \beta \)-dicarbonyl \( \alpha \)-\chem{CH_{2,3}} donors while excluding carboxylic acids and
  acyl chlorides.
\end{remark*}

\section{Order Invariance and Bounded Projection}\label{sec:projection}

\begin{theorem}[Reactant-Order Invariance]\label{thm:commutativity}
  For any \( a, b \in \Z_K \) and \( c \in \Conds_* \), define the canonical pair
  \( \operatorname{canon}(a, b) = (\min(a, b), \max(a, b)) \). Evaluating both slot
  orientations over \( \operatorname{canon}(a, b) \) and selecting the candidate minimizing
  \( (\operatorname{priority}(r_k), \operatorname{SMILES}(P)) \) satisfies
  \begin{equation}\label{eq:commutativity}
    \operatorname{eval}(a, b, c) = \operatorname{eval}(b, a, c).
  \end{equation}
  When no candidate matches under \( c \), no random or fallback substitute is emitted.
\end{theorem}
\begin{proof}
  Both \( \min \) and \( \max \) are symmetric on \( \Z_K \), so
  \( (\ell, u) \coloneqq \operatorname{canon}(a, b) = \operatorname{canon}(b, a) \) for all
  \( a, b \in \Z_K \). Evaluation inspects the fixed finite sequence of transforms
  \( r_k \in \Reacts \) permitted by \( c \) against the identical ordered pair list
  \( [(\ell, u)] \) when \( \ell = u \) and \( [(\ell, u), (u, \ell)] \) when
  \( \ell \neq u \), generating an identical candidate set \( \mathcal{S}(a, b, c) \). If
  \( \mathcal{S}(a, b, c) = \emptyset \), both evaluations return the compatibility failure record
  without substitution. Otherwise, template priorities \( 1, \dots, 8 \) are pairwise distinct, and
  canonical isomeric SMILES serialization assigns a unique string \( \operatorname{SMILES}(P) \)
  to each molecular isomorphism class, so the lexicographic order on the sort key
  \( (\operatorname{priority}(r_k), \operatorname{SMILES}(P)) \) uniquely identifies the winning
  template and canonical product independently of input ordering, giving~\eqref{eq:commutativity}.
\end{proof}

\begin{theorem}[Bounded Projection and Reactant Avoidance]\label{thm:projection}
  Given the standard InChIKey~\( \kappa \), let
  \( h_b = \FNV(\text{\texttt{base:}} \Vert \kappa) \bmod M_b \) and
  \( h_o = \FNV(\text{\texttt{overlay:}} \Vert \kappa) \bmod M_o \).
  For an exclusion list \( E = [e_1, \dots, e_m] \) with \( m \le 2 \) and
  \( e_j \in \Z_{M_b} \), define \( \pi_b(h_b, M_b, E) \) to be \( h_b \) when
  \( M_b < 3 \), and otherwise the result of at most \( m \) unit shifts modulo \( M_b \)
  stopping at the first index not in \( E \). Then:
  \begin{enumerate}
    \item \( 0 \le \pi_b(h_b, M_b, E) < M_b \) and \( 0 \le h_o < M_o \) for all
      \( M_b \ge 2 \), \( M_o \ge 2 \);
    \item whenever \( M_b \ge 3 \), \( \pi_b(h_b, M_b, E) \notin E \) and at most
      two modular increments occur.
  \end{enumerate}
\end{theorem}
\begin{proof}
  For part~(1), reducing modulo \( M_b \ge 2 \) and \( M_o \ge 2 \) directly yields
  \( h_b \in \Z_{M_b} \) and \( h_o \in \Z_{M_o} \). Because every shift step reduces modulo
  \( M_b \), the final value \( \pi_b(h_b, M_b, E) \) always lies in \( \Z_{M_b} \).

  For part~(2), assume \( M_b \ge 3 \) and consider the shifted sequence
  \( x_t = (h_b + t) \bmod M_b \) for \( t \in \{0, 1, 2\} \). Since \( M_b \ge 3 \), the three
  residues \( x_0, x_1, x_2 \) are pairwise distinct in \( \Z_{M_b} \). Because
  \( \lvert E \rvert \le m \le 2 < 3 \), by the pigeonhole principle at least one index
  \( t^* \in \{0, 1, \dots, m\} \) satisfies \( x_{t^*} \notin E \) (if \( m = 0 \) then
  \( E = \emptyset \) and \( t^* = 0 \); if \( m = 1 \) then \( x_0 \neq x_1 \) guarantees
  \( \{x_0, x_1\} \not\subseteq E \); if \( m = 2 \) then \( \{x_0, x_1, x_2\} \not\subseteq E \)).
  Stopping at the least such \( t^* \) yields \( \pi_b(h_b, M_b, E) = x_{t^*} \notin E \) after
  at most \( 2 \) modular increments.
\end{proof}

\section{Exhaustive Atlas Certification}\label{sec:certificate}

\begin{proposition}[Atlas Reactivity Matrix]\label{prop:chem-certificate}
  Evaluating all \( 32 \times 32 \times 5 = 5{,}120 \) ordered triples
  \( (a, b, c) \in {\Z_K}^2 \times \Conds_* \) yields exactly \( 332 \) reactive triples
  (\( 166 \) under \( \text{auto} \) and \( 166 \) partitioned disjointly across
  \( \Conds \)) and \( 4{,}788 \) non-reactive triples, with \( 0 \) sanitization failures,
  \( 0 \) stoichiometric violations, and \( 0 \) commutativity violations. Every template in
  \( \Reacts \) matches at least eight ordered pairs (\cref{tab:reactions}). Under the pinned
  two-dimensional vector-to-PNG renderer, every \( 320 \times 120 \) skeletal diagram encodes a
  valid PNG of at most \( 8{,}743 \) bytes, strictly below the \( 131{,}072 \)-byte limit.
\end{proposition}
\begin{proof}
  Exhaustive enumeration across all \( 32^2 = 1{,}024 \) substrate pairs and all \( 5 \) regimes
  in \( \Conds_* \) computes every transform match, sanitization check,~\eqref{eq:stoichiometry},
  symmetry check~\eqref{eq:commutativity}, and two-dimensional diagram encoding in finite time.
  It also verifies directly that zero ordered pairs in \( \Atlas^2 \) react under more than one
  explicit condition in \( \Conds \); because \( \Reacts \) partitions the active conditions
  disjointly across \( \text{heat} \), \( \text{acid} \), \( \text{base} \), and
  \( \text{palladium} \), every compatible pair therefore attains the identical winning
  transform under both \( \text{auto} \) and its unique owning condition. Summing the explicit and
  automatic columns of~\cref{tab:reactions} gives \( 166 + 166 = 332 \) reactive triples and
  \( 3{,}930 + 858 = 4{,}788 \) non-reactive triples with zero failures.
\end{proof}

\begin{table}[t]
  \centering
  \caption{Exhaustive reaction counts across the \( 32 \times 32 \) atlas and \( \Conds_* \).\label{tab:reactions}}
  \begin{tabular}{@{}llrr@{}}
    \toprule
    Template & Condition & Explicit pairs & Auto pairs \\
    \midrule
    Diels--Alder \([4 + 2]\) cycloaddition & heat & 32 & 32 \\
    Fischer esterification & acid & 18 & 18 \\
    Imine condensation & acid & 36 & 36 \\
    Acylation / amide coupling & base & 16 & 16 \\
    Williamson \( \SNtwo \) substitution & base & 12 & 12 \\
    Aldol condensation & base & 32 & 32 \\
    Michael conjugate addition & base & 12 & 12 \\
    Suzuki--Miyaura cross-coupling & palladium & 8 & 8 \\
    \midrule
    Total reactive & & 166 & 166 \\
    Total non-reactive & & 3{,}930 & 858 \\
    \bottomrule
  \end{tabular}
\end{table}

\section{Chemical Scope and Limitations}\label{sec:limitations}

To keep the specification faithful to real organic chemistry without claiming predictive physics it
does not implement, the following boundaries apply explicitly:

\begin{enumerate}
  \item \textbf{Rule-based constitutional level.} Reactivity is governed solely by curated 2D
    SMARTS substructure matching and valence sanitization, not by quantum-chemical activation
    barriers, thermodynamic equilibrium constants, solvent effects, or solution \(\mathrm{p}K_{\mathrm{a}}\).
    In particular, aldol template \( r_6 \) accepts aldehyde, ketone, enone, and Knoevenagel
    \( \beta \)-dicarbonyl \( \alpha \)-\chem{CH_{2,3}} donors (including crotonaldehyde-type
    self-condensation of aldehydes) while explicitly excluding carboxylic acids and acyl chlorides.
  \item \textbf{Deterministic tie-breaking for selectivity.} Asymmetric dienes, unsymmetrical
    ketones, or multiple reactive sites can admit constitutional isomers or unannotated stereocenters
    (\emph{endo}/\emph{exo}, \(E\)/\(Z\), \(R\)/\(S\));~\cref{thm:commutativity} resolves multiple
    sanitized candidates strictly by transform priority and lexicographic canonical SMILES rather
    than kinetic branching ratios.
  \item \textbf{Single-step bimolecular stoichiometry.} No iterative polymerization, protecting-group
    manipulation, or explicit counterion solvation is modeled; stoichiometric conservation holds at
    the level of neutral molecular reactant graphs plus declared small-molecule byproducts.
\end{enumerate}

\section{Conclusion}\label{sec:conclusion}

Restricting synthesis evaluation to a curated \( 32 \)-substrate organic atlas and eight balanced
bimolecular reaction templates delivers strict heavy-atom stoichiometry, deterministic order
invariance, and pigeonhole-guaranteed reactant-avoiding catalog projection whenever
\( M_b \ge 3 \), while surfacing honest incompatibility diagnostics whenever two functional
groups fail to match.

\begin{thebibliography}{13}

  \bibitem{Baumann1886}
  E.~Baumann,
  Ueber eine einfache Methode der Darstellung von Benzo\"es\"aure\"athern,
  \emph{Ber.\ Dtsch.\ Chem.\ Ges.} \textbf{19} (1886), 3218--3222.

  \bibitem{deMouraUllrich2021}
  L.~de~Moura and S.~Ullrich,
  The Lean~4 theorem prover and programming language, \emph{CADE~28}, LNCS \textbf{12699}
  (2021), 625--635.

  \bibitem{DielsAlder1928}
  O.~Diels and K.~Alder,
  Synthesen in der hydroaromatischen Reihe, \emph{Justus Liebigs Ann.\ Chem.} \textbf{460}
  (1928), 98--122.

  \bibitem{FischerSpeier1895}
  E.~Fischer and A.~Speier,
  Darstellung der Ester, \emph{Ber.\ Dtsch.\ Chem.\ Ges.} \textbf{28} (1895), 3252--3258.

  \bibitem{Heller2015}
  S.~R.~Heller, A.~McNaught, I.~Pletnev, S.~Stein, and D.~Tchekhovskoi,
  InChI, the IUPAC International Chemical Identifier, \emph{J.\ Cheminform.} \textbf{7} (2015), 23.

  \bibitem{Mathlib2020}
  The mathlib Community,
  The Lean mathematical library, \emph{CPP~2020}, ACM (2020), 367--381.

  \bibitem{Michael1887}
  A.~Michael,
  Ueber die Addition von Natriumacetessig- und Natriummalons\"aure\"athern zu den Aethern
  unges\"attigter S\"auren, \emph{J.\ Prakt.\ Chem.} \textbf{35} (1887), 349--356.

  \bibitem{MiyauraSuzuki1995}
  N.~Miyaura and A.~Suzuki,
  Palladium-catalyzed cross-coupling reactions of organoboron compounds, \emph{Chem.\ Rev.}
  \textbf{95} (1995), 2457--2483.

  \bibitem{Schiff1864}
  H.~Schiff,
  Mittheilungen aus dem Universit\"atslaboratorium in Pisa: eine neue Reihe organischer Basen,
  \emph{Justus Liebigs Ann.\ Chem.} \textbf{131} (1864), 118--119.

  \bibitem{Schotten1884}
  C.~Schotten,
  Ueber die Oxydation des Piperidins,
  \emph{Ber.\ Dtsch.\ Chem.\ Ges.} \textbf{17} (1884), 2544--2547.

  \bibitem{Weininger1988}
  D.~Weininger,
  SMILES, a chemical language and information system.\ 1.\ Introduction to methodology and encoding
  rules, \emph{J.\ Chem.\ Inf.\ Comput.\ Sci.} \textbf{28} (1988), 31--36.

  \bibitem{Williamson1850}
  A.~W.~Williamson,
  Theory of \ae{}therification, \emph{Philos.\ Mag.} \textbf{37} (1850), 350--356.

  \bibitem{Wurtz1872}
  A.~Wurtz,
  Sur un ald\'ehyde-alcool,
  \emph{Bull.\ Soc.\ Chim.\ Fr.} \textbf{17} (1872), 436--442.

\end{thebibliography}

\end{document}
